\section{模型假设}

为确保模型的合理性和可操作性，本文基于现有技术发展趋势和工程实践经验，提出以下三条核心假设：

\begin{enumerate}
    \item \textbf{太空电梯系统于2050年达到设计运力并稳定运营。}
    
    本假设的核心内容是：三座银河港在2050年任务启动时已全面投入运营，每座港口提供\textbf{17.9万公吨}的年运力。这一假设基于MCM机构公布的太空电梯建设时间表，该时间表考虑了石墨烯系绳材料的工业化生产、地球港口基础设施建设、顶点锚点部署以及升降机系统调试等关键里程碑。
    
    \textit{合理性依据}：当前碳纳米管和石墨烯材料的拉伸强度已接近理论预测值的50\%以上，预计到2040年可实现太空电梯所需的强度-重量比。日本和中国的研究团队已完成小规模原型测试，验证了电磁推进升降机的技术可行性。国际太空电梯联盟（ISEC）的路线图预测2045年前完成首座银河港建设。
    
    \textit{对模型的影响}：该假设意味着模型未考虑建设爬坡期，即假设系统从任务第一天起即以满负荷运转。实际部署可能需要2-5年的产能爬坡，期间运力逐步从50\%提升至100\%。这可能导致模型低估初期运输进度。
    
    \textit{缓解措施}：敏感性分析中，我们检验了初始运力降低20-40\%的情景，结果显示对总工期影响控制在8\%以内，表明模型对此假设具有一定鲁棒性。
    
    \item \textbf{火箭技术在2050年达到预测性能水平，实现可重复使用和成本大幅降低。}
    
    本假设的核心内容是：先进猎鹰重型或同级别火箭在2050年可实现单次\textbf{100-150公吨}月球运载能力（模型取125公吨均值），运输成本降至\textbf{1500美元/公斤}。这一成本较当前水平（约2500-3000美元/公斤）下降40-50\%。
    
    \textit{合理性依据}：SpaceX的可重复使用技术已将地球轨道发射成本从2010年的约10,000美元/公斤降至2024年的约2,700美元/公斤，降幅超过70\%。按照学习曲线理论，每累计发射量翻倍，成本下降15-20\%。2050年前全球累计发射量预计增长10倍以上，支撑40-50\%的进一步成本下降。NASA的Artemis计划和中国的载人登月计划将推动大型运载火箭技术的持续进步。
    
    \textit{对模型的影响}：若技术进步慢于预期，实际成本可能高于假设值，导致纯火箭方案和混合方案的总成本上升。反之，若出现颠覆性技术突破（如核热推进商业化），成本可能更低。
    
    \textit{缓解措施}：模型采用保守的中间估计值（1500美元/公斤），并在敏感性分析中检验了成本上浮50\%和下降30\%的情景，确保结论的稳健性。
    
    \item \textbf{运营环境允许两系统按设计参数持续稳定运行，不发生灾难性系统损失。}
    
    本假设的核心内容是：太空电梯不因空间碎片撞击发生系绳断裂，火箭发射不因极端天气或技术故障出现连续失败潮。常规故障（电梯\textbf{2\%}年停机率、火箭\textbf{3\%}单次失败率）已纳入模型，但不考虑需要多年修复的灾难性事件。
    
    \textit{合理性依据}：系绳设计采用多股冗余结构，单根损坏不导致系统失效。太空碎片监测和规避技术到2050年将更加成熟，可提前预警并调整升降机运行时间窗口。火箭发射的历史失败率从1960年代的约15\%降至2020年代的约3\%，技术成熟度持续提升。
    
    \textit{对模型的影响}：该假设使模型聚焦于正常运营条件下的优化问题，简化了分析复杂度。若发生灾难性事件（如系绳断裂需5年重建），实际工期将显著延长。
    
    \textit{缓解措施}：建议MCM机构在混合方案中保持一定比例的火箭运力作为战略冗余，以应对电梯系统可能的长期故障。本文的最优55.9\%电梯配置已隐含此冗余考量。
\end{enumerate}
